% PDF 73--78；原书第 43--48 页。
\section{数据}
\label{sec:ch02-data}
前面主要讨论了大语言模型的技术概念。除了模型架构本身，要理解大语言模型如何工作，训练数据也同样重要。本节介绍常见的训练数据类型及其对大语言模型能力的影响。

\subsection{语言模型的预训练数据集}
自 2018 年以来，迁移学习成为自然语言处理各领域的主流。那一年，先后发布了 ULMFiT、GPT 和 BERT 三个重要的语言模型。三者的架构差异很大，但有一项共同做法，即只使用自监督语言建模目标（self-supervised language modeling objective）进行预训练，再用特定任务的有标签数据微调。由此可以借助大量文本获得一般的语言理解能力，无需为文本标注，而有标签数据往往难以取得。本节介绍语言模型预训练最常用的数据来源。

预训练的目标是让大语言模型具备通用的语言理解能力和世界知识。因此，选用的训练数据应覆盖广泛的主题，使用丰富的词汇，并包含具有代表性的书面语言模式分布。数据当然还必须足够多。常用来源如下。
\begin{itemize}
\item \textbf{网页抓取（web scraping）。}从某个网页中的链接出发自动收集网页，再沿新收集网页中的链接继续收集，依此类推。这类数据语言范围广，但质量参差不齐。互联网上的俚语、拼写错误和其他不规范表达，有助于模型适应不同输入。另一方面，许多文本可能无法理解或与事实相悖，若清理不充分，会给模型带来不利影响。Common Crawl 是最知名的公开网页抓取数据。
\item \textbf{维基百科。}维基百科数据有多项好处。首先，它提供大量事实信息，通常编辑得较好，格式也统一，较少出现普通网页中的数据质量问题。维基百科还收录多种语言的文章，便于将训练范围扩展到英语以外。
\item \textbf{图书。}小说提供了优质的叙事文本，描述人们如何思考，如何与他人及环境互动。这类语言在维基百科等知识库中找不到，后者只用第三人称叙述事件。大多数图书也很适合用来建模长程依赖。显然，故事书中的许多信息是虚构的。
\item \textbf{代码。}随着生成模型能力增强，代码生成成为常见应用。人们经常利用 GitHub 和 StackExchange 的数据训练能生成代码的模型。值得注意的是，使用代码训练也可能增强大语言模型在其他逻辑推理任务中的能力（\refcite{ch02-fu}{Fu and Khot, 2022}）。
\end{itemize}

按早期 Transformer 模型的训练规模，通常选用上述来源中的一两种。到了现代大语言模型的规模，更常见的做法是把这些来源乃至更多来源合在一起，利用各自的长处。The Pile（\refcite{ch02-gao}{Gao et al., 2020}）汇集了 22 个来源的语料，包括法律文本、医学文本、学术论文以及 GitHub 上的代码。研究表明，与在 Common Crawl 等来源较单一的语料上训练的模型相比，使用这些数据可以改善下游任务表现。ROOTS 语料库（\refcite{ch02-laurencon}{Laurençon et al., 2023}）把这一思路进一步扩展，从数百个来源纳入 46 种自然语言和 13 种编程语言。

\begin{table}[!htbp]
\centering
\caption{大语言模型预训练中广泛采用的若干语料库。}
\label{tab:ch02-corpora}
\small
\begin{tabularx}{\linewidth}{@{}lX@{}}
\toprule
语料库 & 来源\\
\midrule
BookCorpus & 图书\\
Wikitext103 & 维基百科\\
Common Crawl & 互联网\\
OpenWebText & 互联网\\
The Pile & 互联网、学术研究、图书、对话、代码\\
ROOTS & 高资源和低资源语言、互联网、代码\\
\bottomrule
\end{tabularx}
\end{table}

\subsubsection{多语言语料库和平行语料库}
很多大语言模型完全或主要用一种语言训练。要让模型理解多种语言并在其间翻译，需要涵盖所需各种语言的数据。这类数据大致分为两类。
\begin{itemize}
\item \textbf{平行语料库（parallel corpus）。}每个文本样例都有一个对应的第二种语言的译文。训练使用这些语言对，把一种语言作为输入，另一种作为目标。模型预测与目标越接近，得分越高。
\item \textbf{多语言语料库（multilingual corpus）。}包含多种语言的数据，各语言之间没有明确的翻译对应关系。此类语料适用于语言建模目标，无法直接用于上述平行语料库所用的机器翻译目标。
\end{itemize}
近年来，现代大语言模型的规模已经足以让模型在少样本设置下表现出良好的翻译能力，无需特意用平行数据训练（\refcite{ch02-workshop}{Workshop et al., 2023}）。翻译能力来自模型对多种语言的共同学习，而无须从明确成对的不同语言文本中学习。

\subsection{数据预处理}
预训练所用的语料库远超人工可审阅的规模。因此，人们采用多种方法过滤可能妨碍模型表现或造成非预期影响的数据。过于偏离语言分布的文本，以及含有冒犯性内容或敏感个人信息的文本，都应剔除。

\begin{figure}[!htbp]
\centering
% 重绘自原书第 45 页的图 2.7。
\begin{tikzpicture}[
  font=\footnotesize\sffamily,
  sheet/.style={draw=black!37,fill=black!9,minimum width=14mm,minimum height=11mm,align=center},
  filterstage/.style={single arrow,draw=black!38,fill=yellow!27,minimum width=17mm,
    minimum height=8mm,single arrow head extend=1mm,inner sep=1pt,align=center}]
\node[sheet] at (.17,.18) {};
\node[sheet] at (.08,.09) {};
\node[sheet] at (0,0) {原始\\数据};
\node[filterstage] at (1.90,0) {质量过滤};
\node[filterstage] at (3.86,0) {去重};
\node[filterstage] at (5.82,0) {移除\\有害数据};
\node[filterstage] at (7.78,0) {词元切分};
\node[draw=green!42!gray,rounded corners=2.6mm,fill=green!21,
  minimum width=19mm,minimum height=9.5mm,align=center] at (9.74,0) {训练\\大语言模型};
\end{tikzpicture}

\caption{将大规模语料用于大语言模型预训练前，通常要经过的一系列步骤。}
\label{fig:ch02-cleaning}
\end{figure}

\subsubsection{低质量数据}
如图~\ref{fig:ch02-cleaning} 所示，预处理第一阶段关注数据的整体质量。原始语料通常很大，因此可以删掉相当一部分出现不适合训练迹象的数据。这一阶段的预处理可以较为粗略。

大规模语料库中常见的一个质量问题，是出现超出模型预期用途的语言。例如，若模型专门为西班牙语应用训练，出现西班牙语以外的语言会降低训练效率。可以用语言分类模型或基于规则的方法过滤这些数据。

另一项有用的预处理操作，是依据异常文本模式作统计过滤（statistical filtering）。例如，文本频繁出现远长于普通单词的字符串，标点密度很高，或者过长和过短的句子占比过高。这些模式都表明，相应文档或文档集合较难推广到其他情境，训练价值因而较低。

\subsubsection{重复数据}
人们对训练数据重复的影响已有大量讨论。\refcite{ch02-hernandez}{Hernandez 等人（2022）}发现，使用重复数据训练可能产生几种不良后果。作为对照，\refcite{ch02-biderman}{Biderman 等人（2023）}的分析表明，在重复数据上训练对模型既无益也无害。即使在最理想的情形下，用重复数据训练也似乎无法充分利用计算资源。因此，预处理时尽可能去掉重复文本，已成为标准做法。

\subsubsection{有害数据}
上述问题主要涉及只用最适合的数据来提高训练效率。还有一些信息可能不适合让模型习得。例如，如果大语言模型生成的内容中出现从网上抓取的真实姓名和电子邮件地址，就可能带来问题。训练数据中的毒性和偏见也值得关注。应对这些因素更为复杂，后续章节会再讨论，但在预处理阶段就尽可能剔除冒犯性语言，是有价值的。

\subsubsection{文本规范化}
有些数据的问题较轻，需要清理，但无须完全删去文本。例如，从网页抓取的数据自然会有残留的 HTML 标签，应将它们去掉。另一个常用步骤是 Unicode 规范化（Unicode normalization）。等价的字符串可能有多种编码方式，通常可以选用一种方法尽量统一表示，无须让模型自行学习这些等价关系。类似地，也可以按需把所有文本转为小写，使模型不再把大写字母当作不同的字符。

\subsubsection{词元切分}
预处理完成后，数据用于训练第 2.3.3 节所述的一类词元切分器（tokenizer）。这一步必须先于大语言模型训练，因为切分结果是模型的输入。常见做法是使用现有的词元切分器，免去从头训练，但前提是所用数据来源相近。首先，切分方式必须反映训练数据包含的语言，包括编程语言。此外，对话数据可能偏向“thx”“omg”一类缩写词元，而学术文献中表示专业术语的词元，其分布可能很不一样。

数据按文本块送入词元切分器，每块映射为一串词元。为提高效率，词元表示成长度为 $l$ 的整数向量，$l$ 等于子词的数量。模型第一层又称嵌入层（embedding layer），形状为 $n\times m$，其中 $n$ 是词元切分器习得的词元总数，$m$ 是预先设定的嵌入维数。因此，切分器的输出是一列查找索引，按索引为原始输入识别出的每个词元查得一个 $m$ 维向量。文本由此转换为 $l\times m$ 的浮点数矩阵，可以送入模型开始学习。

\subsection{数据对大语言模型的影响}
如前所述，训练大语言模型有多种数据来源。Gopher 的实验结果（\refcite{ch02-rae-2022}{Rae et al., 2022}）表明，调整各来源所占的数据比例，会明显影响大语言模型在多种下游任务上的整体表现。一般来说，数据多样性往往能改善许多任务的表现，同时也要考虑模型的预期用途。构建聊天机器人时，通常希望训练数据中有相当一部分是对话。反过来说，除非聊天机器人用于提供法律建议，否则没有必要纳入大量法律文档。

\begin{keypoint}
模型在预训练中接触的数据量影响很大。Chinchilla 发布时，这一点已十分明确（\refcite{ch02-hoffmann}{Hoffmann et al., 2022}）。它表明此前的大语言模型训练不足。在追求扩大模型规模所带来的能力时，人们误判了数据规模的作用。Chinchilla 研究者通过实验，尝试建立一个公式，给定算力预算后确定最合适的参数数量和训练词元数量。他们发现，模型和数据的规模应大致按比例增加，与过去强调增加参数的做法明显不同。这说明，有很强能力的大语言模型可以比先前所认为的更小。遵循这些原则，可以更有效地利用预训练预算，微调和推理的成本也会下降。
\end{keypoint}

\subsection{特定任务的数据集}
出于研究需要，自然语言处理“任务”常被用作一种通用尺度，近似衡量给定模型在各种真实场景中的表现。大部分特定任务的数据集都经过精心整理和标注，以便用于监督训练和评估。因此，它们往往比大语言模型预训练所用的超大无标签数据集小得多。

特定任务的数据集通常预先分成训练集（training set）和测试集（test set），确保所有研究者在同样的样例上训练和测试。在这种标准化的数据集上评估表现，可以直接比较不同的架构与训练策略。值得注意的是，大语言模型通常能在未见过相应训练数据的样例时，就在测试集上取得不错的评估结果。这称为零样本学习。
